% =====================================================================
% main.tex
% *** SPLIT VERSION -- THIS IS THE CURRENT/LATEST FILE ***
% Split off from paper-draft-v1.tex on 2026-07-21 into an Overleaf-style
% main.tex + refs.bib layout, at user request. paper-draft-v1.tex is left
% in place unmodified as a single-file historical copy; do not edit it
% further -- edit this file instead. The only change relative to
% paper-draft-v1.tex is mechanical: the inline thebibliography block was
% moved out to refs.bib (same \bibitem keys, now BibTeX entries) and
% replaced below with \bibliographystyle + \bibliography. No wording,
% numbers, or claims were touched. Inline citations in the running text
% are still plain "(Author et al., Year)" text, not \cite{} commands --
% the citation-key pass into the body has still not been done (tracked
% in wiki/meta/manuscript-status.md, open issue 4); this split did not
% attempt it.
%
% Converted from wiki/thesis/paper-draft-v1.md (Markdown) on 2026-07-21.
% Pure format conversion -- no wording, numbers, or claims were added,
% removed, or changed relative to the Markdown source. Only Markdown
% syntax was translated into LaTeX syntax (bold/italic/strikethrough,
% tables, code spans, section headings, special characters).
%
% Patched 2026-07-21 (same day, later pass) to re-sync with three .md
% revisions: (c) confidence-vs-PLIF-recovery correlation recomputed
% against hybrid values (rho=0.046, p=0.682, n=81), (d) Phase 2
% non-learned baselines B1/B3 executed (Sec. 3.2, 4.5), (e) PDBBind
% subset data-source substitution disclosed (Sec. 3.2). Same
% patch-not-rewrite rule applied: only these three changes were ported.
%
% Class: standard article (no ICML/other venue template exists yet in
% this project as of 2026-07-21 -- see manuscript-status.md). Swap the
% \documentclass line and re-flow when a target-venue template is
% chosen; section/subsection structure below is deliberately plain so
% that move is low-friction.
%
% NOTE for whoever does the citation-key pass later: inline citations
% in the running text remain plain "(Author et al., Year)" text, not
% \cite{} commands, because no citation-key pass has been done yet
% (see wiki/meta/manuscript-status.md, open issue 4). The bibliography
% is now refs.bib (BibTeX), with the same \bibitem/citation keys as
% before, but nothing in the body text points to them via \cite yet.
% =====================================================================
\documentclass[11pt]{article}

\usepackage[utf8]{inputenc}
\usepackage[T1]{fontenc}
\usepackage[margin=1in]{geometry}
\usepackage{booktabs}
\usepackage{longtable}
\usepackage{array}
\usepackage{url}
\usepackage[normalem]{ulem} % \sout{} for retracted/invalidated numbers

\newcommand{\notebox}[2]{%
  \par\medskip\noindent\fbox{\begin{minipage}{0.94\linewidth}\small\textbf{#1} #2\end{minipage}}\par\medskip
}

\title{Interaction-Aware Rescoring for DiffDock-L under Cross-Docking\\ Distribution Shift\\[0.3em]\large Manuscript Draft v1}
\author{Authors to be added}
\date{Draft version --- July 21, 2026}

\begin{document}
\maketitle

\notebox{Draft status note (2026-07-20).}{This is a from-scratch rewrite, not a restoration of the
earlier \path{paper-draft-v1.md} (deleted at user request in a prior session). Source material:
\path{team-report-2026-07-16} (lit review, experiment design, and an earlier full-draft version
of Introduction--Experiments), \path{phase2-rescoring-design-2026-07-19} (the current, more
concrete Phase 2 protocol), and the measured Phase 1 results confirmed in \path{log} and
\path{hot} (2026-07-19). All numbers below are traceable to
\path{results/dockgen_pilot/master_eval_table_n122_with_plif.csv} and the accompanying
\path{eval_scripts/} in the \path{diffdock-interaction-rescoring} repository. \textbf{Phase 2
(the rescoring model itself) has not been executed} --- every Phase 2 number in this draft is a
design specification or a hypothesis, never a result, and is labeled as such throughout.

Target venue (per \path{team-report-2026-07-16} \S5.2, unchanged): ICML workshop track, 4--6
pages. This draft is longer than that budget and will need trimming (mainly Related Work) once
Phase 2 results exist and the two phases are confirmed to share one submission.}

\notebox{ERRATUM (2026-07-21).}{\textbf{The PLIF interaction-recovery rate and the
confidence-vs-PLIF-recovery correlation originally reported in this draft (pooled recovery
15.47\%, 138/892, n=80; Spearman rho = $-0.149$, n=78) are INVALIDATED. Do not cite these two
numbers.} A code audit (cross-checking \path{eval_scripts/compute_plif.py} --- this project's
``V1'' PLIF pipeline --- against \path{diffdock_dockgen_inputs.csv}) found that V1 scores every
predicted ligand pose's interaction fingerprint against the \textbf{crystal protein file}
(\path{protein_processed.pdb}) rather than the \textbf{ESMFold-predicted protein file DiffDock-L
actually docked into} (\path{holo_aligned_predicted_protein.pdb}). Every complex's PLIF
fingerprint was therefore computed in a structural context DiffDock-L never used, which
invalidates the resulting numbers regardless of how correctly the downstream pooling and
correlation arithmetic were done. Every occurrence of the retracted numbers below is marked
\textbf{[INVALIDATED --- see \S4.4a]} and left in place rather than deleted, to keep an honest
record of what was measured and later retracted. \textbf{Not affected:} top-1 RMSD $<$ 2~\AA{}
success (14/122 = 11.48\%) and ECE (0.203) --- these come from a separate computation path
(\path{build_master_table.py}, \path{compute_ece.py}) that does not depend on which protein file
PLIF scoring used.}

\notebox{Update (2026-07-21, hybrid recomputation complete).}{The manuscript now has a valid
replacement PLIF recovery number. A hybrid pipeline combining V2's correct protein-file handling
(scoring against the ESMFold-predicted protein DiffDock-L actually docked into) with V1's
non-redundant 9-type ProLIF interaction list was run over the matched n=82 subset of DockGen-E.
Rank-1 (confidence-selected) pooled PLIF interaction recovery is \textbf{47/509 = 9.23\%}
(per-complex mean 8.82\%, median 0\%); the rank1--10 oracle (best of the 10 retained candidates
per complex) is \textbf{64/509 = 12.57\%} (per-complex mean 11.86\%, median 0\%), an oracle
headroom of \textbf{+3.34 percentage points pooled (+3.05 points per-complex mean)}. One complex
(\path{2wwc_1_CHT_2}) has zero true interactions under this 9-type list, so its per-complex
recovery is undefined (NaN); it contributes 0/0 to the pooled figures and is excluded from the
mean/median (n=81 non-null). These two numbers (9.23\% baseline, 12.57\% oracle) are this draft's
reported PLIF recovery figures, replacing the retracted 15.47\%. Source:
\path{hybrid_plif_baseline_matched82.csv} and \path{hybrid_plif_oracle_matched82.csv} (scripts
\path{hybrid_plif_baseline.py}, \path{hybrid_plif_oracle.py}), under
\path{results/dockgen_pilot_eval/} on the project's HPC environment. The confidence-vs-PLIF-recovery
correlation against these hybrid values was still unmeasured at this point in the project --- see
the next update below, which reports it. An interaction between the two oracle-headroom results
(RMSD: +6.56 points; PLIF: +3.34 points) also now informs the Phase 2 framing --- see \S3.2 and
\S4.5. Full detail in \S4.4a.}

\notebox{Update (2026-07-21, confidence-vs-PLIF-recovery correlation now computed): no significant
relationship found.}{All 82 hybrid baseline rows matched a \path{confidence_rank1} value; the one
complex with zero true interactions and undefined recovery (\path{2wwc_1_CHT_2}) was excluded,
leaving n=81. Spearman rho = $0.046167845505$, two-sided $p = 0.682340974058$; an independent
average-rank Pearson calculation reproduced the same rho exactly. \textbf{The previous V1-derived
rho = $-0.149$ remains INVALIDATED and was not reused.} As with that earlier, now-retracted result,
we read this only as a limited, non-causal statement: no statistically significant monotonic
relationship between DiffDock-L's confidence score and this corrected PLIF-recovery rate was
observed in this sample. Source: \path{eval_scripts/correlate_hybrid_plif.py},
\path{results/dockgen_pilot_eval/hybrid_plif_confidence_spearman.json}. Full detail in \S4.4a.}

\begin{abstract}
DiffDock-L selects its output pose using a learned confidence score trained to predict whether a
sampled pose falls within 2~\AA{} RMSD of the native structure. This score has been validated
only as a ranking signal on in-distribution (PDBBind) data; no published study has measured its
formal calibration under cross-docking distribution shift, and no study has measured whether it
predicts recovery of the specific protein--ligand interactions (hydrogen bonds, halogen bonds,
$\pi$-stacking) that determine whether a pose is chemically, not just geometrically, correct. We
run DiffDock-L on all 122 complexes of the DockGen-E cross-docking benchmark and evaluate the raw
confidence score directly. Of 122 complexes, 91 produced RMSD-scorable poses; on this subset,
top-1 RMSD $<$ 2~\AA{} success is 14/122 (11.48\%, denominator fixed at 122 with runtime failures
counted as non-success) and expected calibration error (ECE) is 0.203. An earlier reported pooled
PLIF interaction recovery of 15.47\% (138/892, n=80) and its confidence-vs-PLIF-recovery rank
correlation (Spearman rho = $-0.149$, n=78) \textbf{were computed against the wrong (crystal, not
ESMFold-predicted) protein file and are INVALIDATED [see \S4.4a].} A hybrid recomputation using
the correct protein file and a non-redundant 9-type interaction list gives a valid replacement:
pooled PLIF interaction recovery of the confidence-selected rank-1 pose is \textbf{9.23\%}
(47/509, matched n=82), against a rank1--10 oracle upper bound of \textbf{12.57\%} (64/509) --- an
oracle headroom of +3.34 percentage points. A recomputed Spearman correlation between confidence
and this corrected PLIF-recovery rate finds no statistically significant monotonic relationship
(rho = 0.046, two-sided $p = 0.682$, n=81); the previous, now-invalidated rho = $-0.149$ was not
reused. Taken together, the valid results show low top-1 RMSD success, poor calibration, and low
PLIF recovery with only modest room to improve it by reranking alone, and no statistically
significant relationship (at this sample size) between confidence and this chemical-correctness
signal. An oracle-headroom analysis of
DiffDock-L's already-sampled candidate pool helps explain why the two
correctness criteria give such different pictures. RMSD headroom is substantial (14/122 to
22/122, +6.56 percentage points): a better-RMSD pose is often already sitting in the candidate
pool, and confidence simply is not picking it, a selection problem reranking could fix. PLIF
headroom is small (+3.34 points): the candidate pool rarely contains a chemically much better
alternative in the first place, so no amount of reranking recovers much there (\S4.5). Building
on this diagnosis, we describe (design only; not yet executed) an Interaction-Aware Rescoring
model that reranks DiffDock-L's already-sampled candidate poses using complex-invariant ProLIF
interaction-type counts, pharmacophore geometry, and ensemble consistency, trained on PDBBind and
evaluated without retraining any part of DiffDock-L, on the same 122-complex cross-docking set.
The RMSD-side headroom alone is our justification for building this model; we do not claim it
will meaningfully improve PLIF-based chemical correctness, and treat any PLIF-recovery gain it
shows as a secondary, exploratory finding rather than a headline result. We report the measured
diagnosis in full and the correction as a pre-registered protocol, including a mandatory
oracle-headroom gate that must pass before the rescoring model is trained at all.
\end{abstract}

\section{Introduction}

Diffusion-based generative models have reframed blind molecular docking from a regression problem
into a sampling problem over a product manifold of ligand translation, rotation, and torsion
angles. DiffDock, and its larger successor DiffDock-L, pair this generative sampler with a
separately trained confidence model that predicts whether a sampled pose satisfies the
conventional 2~\AA{} RMSD success criterion. In practice this confidence score has become the
default mechanism for pose selection: DiffDock-L returns the pose its confidence model ranks
first, and downstream virtual-screening pipelines use the same score to triage predictions. That
use rests on two assumptions that, to our knowledge, no published study has tested directly:
first, that the confidence estimate stays informative when the model is applied outside the
distribution it was trained on; second, that RMSD-based correctness --- the only target the
confidence model was ever trained to predict --- is itself a sufficient proxy for the properties a
downstream user actually cares about, such as whether a pose reproduces the key hydrogen bonds
and other contacts that determine binding.

Two bodies of evidence expose the first assumption as untested. DiffDock's confidence score has
been validated as a ranking quantity --- via Spearman rank correlation with RMSD and via
selective-prediction accuracy gains --- exclusively on in-distribution PDBBind time-splits.
Separately, a growing set of cross-docking and out-of-distribution benchmarks documents that
DiffDock-L's raw pose accuracy is highly protocol-dependent, with reported top-1 success rates
ranging from roughly 5\% to 55\% depending on dataset, sampling budget, and self- vs. cross-docking
setup. No source in either body of evidence connects the two: nobody has measured whether the
confidence score degrades in a calibrated, predictable way as accuracy collapses under
distribution shift, or whether it instead stays silently overconfident. A third, independent line
of work exposes the second assumption. Errington et al.\ (arXiv:2409.20227) show that RMSD-based
pose correctness carries no measured relationship to whether a pose reproduces the reference
protein--ligand interaction fingerprint (PLIF) --- yet no study we are aware of feeds that
chemically orthogonal signal back into a docking model's own pose-selection mechanism.

This project addresses both gaps with a diagnose-then-correct structure. Phase 1 measures the
calibration of DiffDock-L's raw confidence score under cross-docking shift, using the DockGen-E
benchmark (122 complexes with pockets dissimilar to PDBBind), and reports --- for the first time
on this benchmark --- PLIF interaction recovery alongside RMSD-based top-1 accuracy. An earlier
reported version of the PLIF recovery rate (15.47\%) and a test of whether confidence tracks it
(Spearman rho = $-0.149$, n=78, $t(76) = -1.31$) \textbf{[INVALIDATED --- see \S4.4a]: both were
computed by a pipeline that scored the predicted ligand against the wrong (crystal, not
ESMFold-predicted) protein file.} A hybrid recomputation against the correct protein file now
gives a valid pooled PLIF recovery rate for the confidence-selected pose of 9.23\% (47/509,
matched n=82), against a rank1--10 oracle upper bound of 12.57\% (64/509); a recomputed
correlation between confidence and this corrected recovery rate finds no statistically
significant monotonic relationship (Spearman rho = 0.046, two-sided $p = 0.682$, n=81). The
RMSD-based Phase 1 measurements remain complete and valid and are the primary reported result of
this draft:
on the 91 of 122 complexes with RMSD-scorable output, top-1 RMSD $<$ 2~\AA{} success is 11.48\%
(14/122) and ECE is 0.203. Phase 2, building on that diagnosis, proposes an Interaction-Aware
Rescoring model that reranks DiffDock-L's already-sampled candidate poses using PLIF-derived and
pharmacophore features rather than confidence alone. \textbf{Phase 2's mandatory pre-registered
oracle-headroom gate (Step 0) has now been run} on the retained rank1--10 candidate pool: the
RMSD gate clears the pre-registered +5-percentage-point bar (14/122 to 22/122, +6.56 points),
while the PLIF gate does not (9.23\% to 12.57\%, +3.34 points). We read this as two different
diagnoses, not a contradiction --- the RMSD headroom indicates a \textit{selection} failure that
reranking can address, while the small PLIF headroom indicates a \textit{sampling} limitation
that reranking cannot --- and we proceed to build and evaluate the rescoring model on
RMSD-selection grounds alone, without claiming that it will materially improve PLIF-based
chemical correctness (\S3.2, \S4.5). We report the Phase 1 diagnosis and the Phase 2 gate outcome
as measured results and the Phase 2 model itself as a specified, not-yet-run protocol, and we are
explicit throughout about which numbers are measurements and which are targets.

Contributions: (i) the first reported calibration analysis (ECE) of DiffDock-L's raw confidence
score on a leakage-controlled cross-docking benchmark (DockGen-E, n=122), including an honest
accounting of a substantial fraction of complexes (31/122) for which the pipeline produces no
scorable output at all --- plus a PLIF interaction-recovery rate (9.23\% rank-1, 12.57\% oracle,
after a hybrid recomputation correcting an earlier protein-file bug, \S4.4a) and its recomputed
correlation with confidence (Spearman rho = 0.046, two-sided $p = 0.682$, n=81 --- no statistically
significant monotonic relationship, \S4.4a) --- \textbf{an earlier version of both numbers (15.47\%
recovery, rho = $-0.149$ correlation) [INVALIDATED --- see \S4.4a] was computed against the wrong
(crystal, not ESMFold-predicted) protein file and has been superseded by the hybrid recomputation
above}; (ii) a mandatory pre-registered oracle-headroom gate, now executed, showing that the
RMSD-based candidate pool has meaningful headroom (+6.56 points) while the PLIF-based pool does not
(+3.34 points) --- used here to justify Phase 2 on RMSD-selection grounds while explicitly not
claiming a PLIF-correctness improvement --- together with two executed non-learned reranking
comparators (B1 random, B3 classical affinity) that all cluster within an 11.12--12.30\% band
indistinguishable from raw confidence (\S3.2, \S4.5), plus a fully specified, not-yet-executed
Interaction-Aware Rescoring protocol that reranks DiffDock-L's existing candidate poses using
complex-invariant interaction and pharmacophore features; and (iii) an explicit statement of what
remains unmeasured --- stratified (pocket-similarity) calibration, an in-distribution ECE anchor,
and every Phase 2 learned-model (M1/M1-abl) accuracy number --- so that this draft does not
overstate what has actually been shown.

\section{Related Work}

\textbf{Diffusion docking and its confidence model.} DiffDock (Corso et al., ICLR 2023) recast
blind docking as generative sampling over translation, rotation, and torsion, pairing an
SE(3)-equivariant score network with a separately trained confidence model that predicts whether a
pose meets the 2~\AA{} RMSD criterion. On in-distribution PDBBind time-splits this confidence
score attains Spearman correlation 0.68 with RMSD and lifts selective accuracy from 38\% to 83\%
under aggressive abstention. The same source paper (Corso et al., arXiv:2402.18396, which also
introduced DockGen) shows that Confidence Bootstrapping applied to \textbf{DiffDock-S} (not
DiffDock-L) raises DockGen-clusters accuracy from \textbf{9.8\% to 24.0\%} (Table 1, Fig.~4-D).
This claim was independently verified against the arXiv:2402.18396 source PDF (PMC10925391) on
2026-07-22; the earlier draft had misattributed the improved model as DiffDock-L and misreported
the figures as roughly 7\% to 23\%, both now corrected. DiffDock-PP (Ketata et al., ICLR 2023 workshop, arXiv:2304.03889)
extends the same top-1-plus-confidence paradigm to protein--protein docking and reports a large gap
between oracle-selected and confidence-selected poses. In every case we are aware of, the
confidence estimate is validated as a ranking or selective-prediction signal, never through a
formal calibration metric such as ECE, and never re-examined once a richer feature set ---
interaction fingerprints, pharmacophores --- is available at inference time.

\textbf{Cross-docking and out-of-distribution generalization.} A second line of work documents the
fragility of deep-learning docking under distribution shift but measures pose accuracy, not
confidence. PoseX (Jiang et al., ICLR 2026, arXiv:2505.01700), a leakage-controlled self- and
cross-docking benchmark, reports a negative correlation ($r \approx -0.50$) between DiffDock-L's
pocket similarity to training data and its pose RMSD, but evaluates only top-1 RMSD and does not
assess the confidence score itself. FlowDock (Morehead \& Cheng, arXiv:2412.10966) reports
DiffDock-L accuracy collapsing to 4.8\% on DockGen-E under its own evaluation protocol. A
comparative study on SARS-CoV-2 and MERS-CoV main proteases (Liu et al., J. Chem. Inf. Model.
2026) finds vanilla DiffDock confidence-selected top-1 accuracy of 12.3\%, below classical Glide,
on novel viral protease targets. A large-scale LIT-PCBA virtual-screening evaluation (Abo-Dahab et
al., UCSF capstone 2026, arXiv:2605.01681) finds DiffDock confidence performing at or below
classical AutoDock in a single-pose screening setting. PoseBench (Morehead et al., Nature Machine Intelligence 8, 32--41,
2026) ties these degradations to a mechanism, reporting a strong negative correlation between
deep-learning accuracy and training-set similarity --- consistent with memorization --- that is
absent for classical, non-learned docking. Reported DiffDock-L success rates therefore vary by
roughly an order of magnitude ($\approx$5--55\%) with benchmark protocol; any calibration or
rescoring result must fix and report its exact protocol to be comparable to prior numbers at all,
and we do so in \S4.

\textbf{Calibration, confidence, and interaction recovery in adjacent settings.} Formal
calibration analysis exists in this literature mainly for co-folding models, not pose-sampling
docking models. AlphaFold3 (Abramson et al., Nature 630, 493--500, 2024) reports the only
reliability-diagram-style analysis we found in this corpus, and it varies sharply by interaction
category --- a caution against treating a single aggregate calibration number as representative.
ACER (ICML 2026 workshop submission) reports that co-folding confidence correlates weakly with
success on out-of-distribution and allosteric targets. An agentic LLM-based pose reranker (Khiari
et al., 2026) finds self-reported verbalized confidence inversely correlated with accuracy in its
small evaluation. Most directly motivating this work, Errington et al.\ (arXiv:2409.20227) show
that RMSD-based pose correctness --- the exact target DiffDock-L's confidence model is trained on
--- carries no measured relationship to whether a pose reproduces the reference protein--ligand
interactions computed via ProLIF. We are not aware of prior work that closes this loop by feeding
interaction-fingerprint or pharmacophore information back into a trained reranking mechanism for
diffusion-docking poses; existing pose-rescoring work --- GNINA's CNN scoring (Buccheri \&
Rescifina, Molecules 2025), MC-GNNAS-Dock's multi-criteria selection (Cao et al.,
arXiv:2509.26377) --- reranks with classical or graph-based scoring functions, not by augmenting a
diffusion model's own learned confidence output with chemically orthogonal features. The closest
adjacent precedent we found is DiffDock-NMDN (Xia, Gu \& Zhang, \emph{J. Chem. Inf. Model.} 65(3),
1101--1114, 2025), which does rerank DiffDock's sampled poses with a separately trained model, but
that model (NMDN) learns a continuous protein-residue--ligand-atom distance-likelihood potential
via a mixture density network over ESM-2/sPhysNet embeddings --- purely geometric supervision,
with no discrete interaction-type label (hydrogen bond, hydrophobic contact, $\pi$-stacking) or
pharmacophore directional constraint anywhere in its inputs or loss. This is architecturally
distinct from the interaction-fingerprint/pharmacophore features this work proposes to feed back
into reranking, and an independent large-scale evaluation of this exact pipeline (Abo-Dahab et
al., arXiv:2605.01681) reports that DiffDock-NMDN underperforms the unscored DiffDock baseline on
LIT-PCBA (median EF1\% 0.67 vs.\ 1.17), which we read as evidence that a purely distance-based
rescorer does not obviously substitute for the discrete interaction signal this work targets.

\textbf{Positioning.} The confidence-related evidence in this literature splits into three parts
that do not overlap: in-distribution ranking evidence for DiffDock-family confidence, out-of-
distribution accuracy-degradation evidence that never measures confidence, and
interaction-recovery evidence with no link to confidence or rescoring. This work first measures
the missing link between the first two (calibration under cross-docking shift, reported below),
then specifies --- as a not-yet-executed follow-on --- the missing link to the third (a rescoring
model that explicitly uses the interaction signal raw confidence ignores).

\section{Method}

\subsection{Phase 1 --- calibration diagnosis (executed on DockGen-E; broader protocol only partly
run)}

The object under evaluation is the calibration of DiffDock-L's own confidence output, not its raw
accuracy. For each complex, DiffDock-L samples N candidate poses and emits a confidence logit per
pose, trained under a binary cross-entropy-with-logits objective to predict RMSD $<$ 2~\AA{}. We
take sigmoid(logit) of the top-ranked pose as a probability estimate P(RMSD $<$ 2~\AA{}),
consistent with that training objective, and evaluate it against two ground-truth labels: the
conventional RMSD $<$ 2~\AA{} criterion (primary), and ProLIF-computed PLIF interaction recovery
against the reference structure (secondary, chemically orthogonal).

The full designed protocol specifies five comparators: DiffDock-L on PDBBind (in-distribution
anchor), base DiffDock on the same cross-docking complexes (within-family contrast), FlowDock's
plDDT-style confidence (deep-generative comparator), classical docking scores (AutoDock Vina
primary, GNINA CNN score secondary, as non-learned, non-memorization contrasts), and post-hoc
recalibration (temperature/isotonic scaling, to separate a simple scale problem from a directional
failure). \textbf{Of these five, only DiffDock-L on DockGen-E has been run and evaluated as of
this draft} (\S4). The PDBBind anchor, base DiffDock, FlowDock, and the classical-scoring and
recalibration comparators are installed and functional on the project's compute environment (SDSU
HPC ``Innovator'' cluster) but have not yet been executed on this benchmark; we report this
explicitly rather than assume their outcome.

Metrics are ECE, adaptive-binning ECE, Brier score, and reliability diagrams --- planned aggregate
and stratified by pocket similarity to training data (max TM-score for PoseX-CD, PLINDER SuCOS for
DockGen-E) --- plus risk--coverage / AURC as a comparator-semantics-free cross-model metric, and
the confidence-vs-PLIF-recovery relationship as the second, chemically orthogonal evaluation axis.
\textbf{The aggregate (non-stratified) ECE has been computed and remains valid} (\S4). An
aggregate PLIF interaction-recovery rate and its correlation with confidence were also computed
early on (pooled recovery 15.47\%; Spearman rho = $-0.149$, n=78, $t(76) = -1.31$), but
\textbf{[INVALIDATED --- see \S4.4a]: both were computed against the wrong (crystal, not
ESMFold-predicted) protein file.} A hybrid recomputation against the correct protein file, using a
non-redundant 9-type ProLIF interaction list, now gives a valid pooled PLIF recovery rate: 9.23\%
(47/509, matched n=82) for the confidence-selected rank-1 pose against a rank1--10 oracle upper
bound of 12.57\% (64/509, \S4.4a). The confidence-vs-PLIF-recovery correlation has since been
recomputed against this corrected data and finds no statistically significant monotonic
relationship (Spearman rho = 0.046, two-sided $p = 0.682$, n=81; \S4.4a). Stratified reliability
diagrams, AURC, and the classical/recalibration comparators remain designed but unexecuted.

\subsection{Phase 2 --- Interaction-Aware Rescoring (oracle gate and non-learned comparators
B1/B3 executed; learned rescorer M1/M1-abl design only, not trained)}

Phase 2 is defined as a reranking problem, not a pose-generation problem: DiffDock-L already
samples and confidence-ranks multiple candidate poses per complex, and the project's HPC runs
retain the top-10 ranked poses per complex from the Phase 1 inference. Rather than resample, Phase
2 reranks within this existing candidate pool of 10 poses per complex, which requires no
additional GPU inference for the evaluation benchmark and keeps the comparison to raw confidence
on an identical candidate set.

\textbf{Step 0 --- oracle-headroom gate (mandatory, run before any model is trained).} For every
complex, we compute per-rank RMSD and per-rank PLIF recovery across all 10 candidates and compare
three quantities, all on the same 122-complex denominator: (B0) the existing raw-confidence
baseline --- rank-1 top-1 success; (B2) an oracle upper bound --- success if any of the 10
candidates clears the criterion, defining the best any reranking method could possibly achieve;
and (B1) a random-choice lower sanity bound. The pre-registered decision rule is: if oracle
$\approx$ baseline (headroom below roughly 5 percentage points, i.e., fewer than about 6
complexes), rescoring is abandoned as a direction for that criterion, because the correct pose is
not present in the candidate pool for reranking to find --- a limitation to report, not a
rescoring failure. If headroom clears that bar, rescoring proceeds on that criterion. \textbf{This
gate has now been run for both RMSD and PLIF, and the two criteria disagree.} RMSD: baseline
14/122 (11.48\%) to oracle 22/122 (18.03\%), headroom +6.56 percentage points --- clears the bar,
\textbf{PASS}. PLIF (hybrid-recomputed, \S4.4a): baseline 9.23\% (47/509) to oracle 12.57\%
(64/509), headroom +3.34 percentage points --- below the bar, \textbf{FAIL}. We read this pair of
results as measuring two different failure modes, not as a contradiction. The RMSD headroom shows
a \textit{selection} problem: a better-RMSD pose already exists in the confidence-ranked candidate
pool for a meaningful number of complexes, and DiffDock-L's confidence score is simply not
choosing it --- exactly the failure a learned rescorer is built to fix. The small PLIF headroom
shows a \textit{sampling} limitation instead: even an oracle that always chose the single
best-PLIF-recovery candidate out of 10 barely improves over confidence's rank-1 choice, which
means the candidate pool itself rarely contains a chemically much-better alternative to rerank
toward. No reranking method --- including the one proposed here --- can recover interactions that
were never sampled in the first place. Applying the decision rule per-criterion (an effective
logical OR across the two gates, since either criterion clearing the bar is sufficient
justification to proceed) means \textbf{Phase 2 proceeds}, justified by the RMSD-side headroom
alone. We do not extend that justification to PLIF: this draft does not claim, and explicitly
disclaims, that Interaction-Aware Rescoring will materially improve PLIF-based chemical
correctness, since the ceiling for that improvement is measured here at +3.34 points regardless of
how good the reranker is. The PLIF-derived features in M1 below (interaction-type counts,
pharmacophore geometry) are accordingly justified only as auxiliary signal that may help
RMSD-based selection --- consistent with Errington et al.'s finding that RMSD-based correctness
and PLIF recovery are not the same axis (\S2) --- not as a mechanism expected to close the small
PLIF gap itself. Any PLIF-recovery improvement M1 shows is reported as a secondary, exploratory
result, not a claimed contribution.

\textbf{Comparators (Step 0 passed, on RMSD-selection grounds).} All comparators rerank within the
same fixed 10-candidate pool and are evaluated on the identical 122-complex denominator used for
Phase 1: B0 (raw confidence, rank-1 --- the existing baseline), B1 (random-in-candidates), B2
(oracle, upper bound, not achievable by any method), B3 (classical rescoring of the 10 candidates
by Vina score, primary, and GNINA CNN score, secondary --- a non-learned, non-PLIF contrast), M1
(the proposed PLIF-aware learned rescorer), and M1-abl (an ablation of M1 with PLIF features
removed, isolating whether any improvement comes specifically from the interaction signal or from
the other features alone).

\textbf{B1 and B3 have since been executed (2026-07-21); M1 and M1-abl remain unexecuted.} B1
(random-in-candidates) is inherently probabilistic, so it is reported as an expected success rate
rather than a single random draw: summing each complex's fraction of the 10 candidates with RMSD
$<$ 2~\AA{} (\path{n_ranks_lt2A}/\path{n_ranks} in \path{rank_oracle_rmsd.csv}) over the 98 of 122
complexes with per-rank RMSD available gives an expected 13.5667/122 = 11.12\%, on the same fixed
122-complex denominator --- 0.36 percentage points below B0's observed 14/122 = 11.48\%, a
difference too small to read as meaningful given B1 is an expectation, not an integer observation.
B3 (classical rescoring) used GNINA 1.3.3 (\path{--score_only --cnn_scoring none}) to score each of
DiffDock-L's retained rank1--10 candidates against the same ESMFold-predicted protein structure
used at inference, selecting the lowest-affinity pose per complex; all 111 complexes with usable
candidates were processed with zero task failures (SLURM array job 12508307), of which 98 are
RMSD-evaluable. B3 succeeds on 15/122 = 12.30\% of complexes, one complex more than B0 (+0.82
percentage points); a paired comparison shows both methods succeeding on the same 11 complexes,
with 3 B0-only and 4 B3-only successes --- only 7 discordant pairs out of 122 --- so this
one-complex difference is not read as a clear improvement. B0, B1, and B3 therefore all fall
within a narrow 11.12--12.30\% band: neither a random reranking nor a classical affinity-based
reranking of the same candidate pool moves the needle materially over raw confidence. This
motivates M1 as a learned alternative, but it also sets the bar M1 must clear by a defensible
margin --- evaluated with the same paired statistics specified below --- for Phase 2 to represent
a genuine improvement rather than a restatement of this band.

\textbf{Features.} Because inference-time candidates have no reference structure, the model cannot
use PLIF \textit{recovery} (which requires a native pose) as an input feature --- only the
absolute interaction fingerprint each candidate pose itself forms. ProLIF's per-residue
interaction bits are not transferable across complexes with different binding-site residues, so
features are aggregated into complex-invariant quantities: counts of each ProLIF interaction type
(hydrogen bond donor/acceptor, hydrophobic, $\pi$-stacking, $\pi$-cation/cation-$\pi$, halogen
bond, ionic, metal coordination) formed by the candidate pose; pharmacophore geometry summaries
(donor/acceptor/aromatic-center distance and alignment statistics); the raw DiffDock-L confidence
logit; and an ensemble-consistency feature counting how many of the other 9 candidates fall within
2~\AA{} RMSD of this one (a near-zero-cost signal independent of PLIF). These feature dimensions
are identical across PDBBind and DockGen-E, enabling cross-complex training and evaluation.

ProLIF fingerprint computation itself is not universally reliable: in Phase 1 it failed outright on
11 of 91 ($\sim$12\%) RMSD-scored complexes (uint32 interaction-index overflow, missing van der
Waals radii for uncommon metal centers, and one residue-identity mismatch; \S4.4). This
failure-rate figure comes from the same V1 pipeline now found to score against the wrong protein
file (\S4.4a); it is plausible but not confirmed that these outright computation crashes are
independent of that bug and will recur unchanged after recomputation. The same failure modes are
expected to recur when computing per-candidate PLIF features for Phase 2, since they arise from
the same ProLIF call applied to structurally similar poses. Phase 2's feature extraction and the
rescorer itself must therefore treat missing PLIF features as a first-class case (e.g., a
missingness indicator plus fallback to the non-PLIF features) rather than assume every candidate
pose yields a usable fingerprint.

\textbf{Training and leakage control.} M1 and M1-abl are gradient-boosted trees (shallow depth,
strong regularization given expected small effective sample size) trained exclusively on a PDBBind
in-distribution split, disjoint from all cross-docking evaluation data. Training requires running
DiffDock-L inference on PDBBind to obtain the same candidate-pose and feature representation used
at evaluation time --- the only additional GPU cost Phase 2 introduces beyond the existing Phase 1
runs, justified now that Step 0 has passed on RMSD-selection grounds. Supervision targets are RMSD
$<$ 2~\AA{} (primary) and PLIF recovery (secondary, exploratory given the small PLIF-oracle
headroom found in Step 0). No component of DiffDock-L itself --- score network or confidence
network --- is retrained; only the reranking step is learned.

\textbf{PDBBind subset data-source caveat.} For tractable inference cost, M1/M1-abl training uses
a 1,500-complex PDBBind subset drawn (seed=42, random sample) from the project's 16,379-ID PDBBind
train list, rather than the full training set. While provisioning it, the data source cited in
DiffDock's own GitHub README (Zenodo record 6408497, EquiBind-processed complexes) was confirmed
deleted (Zenodo API, checked 2025-01-13) --- a pre-existing defect in the upstream DiffDock
repository, not introduced by this project. A substitute, Zenodo record 7014096 (different
uploader, CC-BY-4.0, PDBBind v2020 refined+general set), was adopted; it uses the same complex IDs
and the standard \path{<PDBID>_protein.pdb}/\path{_ligand.sdf}/\path{_pocket.pdb} file layout
DiffDock/PoseBench expects, but was prepared with a different structure-preparation pipeline
(Schr\"odinger Protein Prep Wizard) than the original (EquiBind's own processing), so protonation
states and exact coordinates may differ slightly between the two sources. This has not been
empirically checked (e.g., via a backbone-RMSD spot check between overlapping complex IDs) and is
disclosed here as an open reproducibility caveat on M1/M1-abl's training data, to be resolved
before those results are reported as findings. M1/M1-abl themselves remain unexecuted as of this
draft.

\textbf{Failure exclusions in the training subset.} GPU inference for M1/M1-abl over this
1,500-complex subset is now complete. DiffDock-L's inference path has no fixed random seed (no
\texttt{--seed} argument and no \texttt{torch.manual\_seed} call; the one fixed seed present in
\path{conformer_matching.py} is used only by a training-time function, confirmed unused at
inference), so many inference failures are stochastic and were retried across multiple rounds
before this final accounting. After all retries, code-level fixes (an OpenBabel fallback for
RDKit-unparsable ligands, a conformer-embedding fallback for chirality-sensitive ligands, and a
backbone-atom-selection fix for non-standard-residue receptors), and a full merge-and-recount of
all output directories, \textbf{1,440/1,500 (96.0\%) complexes yielded valid outputs}. The
remaining 60 (4.0\%) were re-classified by reading each complex's individual log section directly
(a whole-log keyword search was found to misclassify some cases and was discarded as unreliable):
37 SVD-ill-conditioned (non-convergent Kabsch alignment), 12 RDKit ligand-parsing failures that
resisted the OpenBabel fallback, 6 TorchScript \path{radius.py} \texttt{RuntimeError} failures
(confirmed fatal, not a benign warning, by inspecting each complex's log for an explicit
\texttt{Failed for 1 complexes} line), 4 receptor-too-large skips (pipeline's 3,000-residue
threshold), and 1 CSV-field \texttt{AttributeError} (a complex-name field parsed as \texttt{float}
rather than \texttt{str}, newly identified this pass) --- $37+12+4+6+1=60$, and $1{,}440+60=1{,}500$.
These 60 complexes are excluded from M1/M1-abl's training data. The RDKit-parsing failures overlap
with a previously characterized list of 236 PDBBind ligands with the same defect (consistent with
PDBBind's own reported $\sim$16\% RDKit-parsing failure rate across the full dataset,
arXiv:2411.01223), a list previously found to concentrate disproportionately in polysaccharide- and
peptide-like ligands rather than being spread uniformly across chemotypes (\S4.6(j)).
\textbf{RMSD self-consistency against the PDBBind-subset ground truth has not yet been
recomputed for this final 1,440-complex sample}; the previously reported 1,043/1,500 (69.5\%)
figure was computed against the earlier, superseded 1,343-complex accounting and is stale --- it
is not restated here, and should not be rescaled by proportion to approximate a 1,440-complex
figure; a fresh computation is required before this sanity check can be reported.

\textbf{Evaluation and statistics.} M1/M1-abl are compared against B0/B1/B2/B3 on the same
122-complex DockGen-E denominator and the same RMSD $<$ 2~\AA{} and PLIF-recovery correctness
definitions used in Phase 1, so Phase 2 numbers are directly comparable to the Phase 1 baseline
reported in \S4. Comparisons are planned to be paired at the complex level (same 122 complexes
across all methods), using McNemar's exact test for paired top-1 correctness plus bootstrap
confidence intervals for the success-rate delta, given the small expected effective sample size
(91 scored complexes, 14 baseline successes) and correspondingly low statistical power. Effect
sizes (the number of complexes gained or lost) are to be reported alongside any p-value, not in
place of it. DockGen-E is treated as a pilot; PoseX-CD (1,312 leakage-controlled complexes) is the
designated confirmatory extension once directionality and effect size are established on
DockGen-E.

\section{Experiments}

\subsection{Datasets}

\textbf{DockGen-E} (122 complexes, curated by PoseBench as a bioassembly-aware refinement of the
original 189-complex DockGen cross-docking generalization benchmark, Corso et al.,
arXiv:2402.18396): complexes with binding pockets dissimilar to PDBBind training data, used here
as the sole cross-docking test set for the measured Phase 1 results below. \textbf{PoseX-CD}
(1,312 leakage-controlled cross-docking complexes, 2022--2025, from CataAI/PoseX,
arXiv:2505.01700) is designated as a planned confirmatory extension and has not been run.
\textbf{PDBBind} serves as the in-distribution anchor for the full designed Phase 1 protocol and
as the exclusive training split for the Phase 2 rescoring model; neither role has been executed
yet.

\subsection{Environment and executed protocol}

DiffDock-L inference ran on the SDSU HPC ``Innovator'' SLURM cluster (\path{gpu} partition), using
PoseBench's official inference and evaluation pipeline (BioinfoMachineLearning/PoseBench, tag
v1.1.0-1-gc5d728d) with N=10 sampled poses per complex. ProLIF 2.2.0 computed PLIF interaction
fingerprints. Base DiffDock, FlowDock, GNINA, and AutoDock Vina are installed and verified
functional in the same environment but have not been run on this benchmark for the present draft.

\subsection{Denominator methodology}

We fix the evaluation denominator at all 122 DockGen-E complexes and count every complex that does
not produce a scorable pose as non-success, rather than restricting the denominator to complexes
with successfully computed output. An earlier plan to report success rate over only the 111
complexes with any usable inference output was rejected during this project as a methodological
error that would inflate the reported success rate; we instead follow the convention, attributed
in the project's working notes to PocketVina (arXiv:2506.20043), of retaining the original sample
size as the denominator even when a method fails outright on some inputs.

Of 122 complexes: \textbf{11 produced no output at all} (runtime failures during inference or
downstream RMSD-alignment, which failed reproducibly across 2 independent retry rounds
(2026-07-18, 2026-07-19) --- traced in this project's execution logs to two distinct causes, 7
numerically ill-conditioned SVD failures during Kabsch alignment and 4 RDKit ligand-conformer or
downstream exception-handling failures); \textbf{20 produced poses but could not be RMSD-scored},
owing to non-standard ligand chemistry that the current DiffDock-L/PoseBench pipeline does not
parse correctly (peptide-like ligands, ligands with repeated glycine residues, and metal-ion
ligands); and \textbf{91 were fully scored}. We flag one documentation discrepancy for the record
and resolve it in favor of the primary source: this project's GitHub repository README attributes
all 11 runtime failures to a single ``group-aggregation bug in PoseBench's fork of
\path{inference.py}'', while the more granular HPC execution log --- which records the exact
error text and complex ID for each of the 11 failures (e.g., \path{linalg.svd: ... ill-conditioned
or has too many repeated singular values} for the 7 SVD cases; \path{'ValueError' object has no
attribute 'RemoveAllConformers'} and \path{inference.py:449: NameError: name 'idx' is not defined}
for two of the 4 RDKit/exception-handling cases) --- attributes them to the two causes above. No
occurrence of ``group-aggregation'' or a matching failure signature appears anywhere in the
execution log's record of these 11 complexes; the log's only group-aggregation-related issue is a
separate, already-fixed bug (trailing \path{_N} suffix stripping in output directory names,
symlink-patched, log entry 2026-07-19) that affected 13 \textit{different}, already-successfully-
scored complexes, not these 11 failures. We therefore treat the execution-log breakdown as the
correct account per this project's evidence-priority order (primary run logs over a summary
README) and flag the README wording as incorrect, to be corrected in a follow-up pass rather than
left as an open discrepancy.

\subsection{Results (measured, 2026-07-19; PLIF figures superseded and correlation recomputed
2026-07-21, see \S4.4a)}

\begin{longtable}{@{}p{4.3cm}p{2.6cm}p{7.3cm}@{}}
\toprule
\textbf{Metric} & \textbf{Value} & \textbf{Denominator / basis} \\
\midrule
\endhead
Top-1 RMSD $<$ 2~\AA{} success rate &
  \textbf{14/122 = 11.48\%} &
  fixed denominator 122 (all complexes; 11 runtime failures and 20 unscored complexes counted as
  non-success) \\[0.4em]
--- reference figure &
  15.38\% &
  91 RMSD-scored complexes only (not the primary denominator) \\[0.4em]
--- reference figure &
  12.61\% &
  111 complexes with any usable inference output (not the primary denominator; rejected as a
  reporting basis, see \S4.3) \\[0.4em]
Expected calibration error (ECE) &
  \textbf{0.203} &
  91 scored complexes; confidence = sigmoid(raw logit) as a P(RMSD $<$ 2~\AA{}) estimate,
  consistent with DiffDock's BCE-with-logits training objective; the 20 unscored complexes are
  excluded from this calculation rather than treated as incorrect \\[0.4em]
PLIF interaction recovery (pooled, rank-1 confidence baseline) --- original computation
  \textbf{[INVALIDATED, \S4.4a]} &
  \sout{15.47\%} (138/892) &
  pooled across 80 of the 91 scored complexes; \textbf{computed against the wrong (crystal, not
  ESMFold-predicted) protein file --- see \S4.4a} \\[0.4em]
PLIF interaction recovery (pooled, rank-1 confidence baseline) --- \textbf{hybrid recomputation,
  valid} &
  \textbf{47/509 = 9.23\%} &
  matched n=82; per-complex mean 8.82\%, median 0\%; correct (ESMFold-predicted) protein file +
  non-redundant 9-type interaction list; one complex (\path{2wwc_1_CHT_2}) has 0 true interactions
  and is undefined (NaN), excluded from mean/median (n=81 non-null), contributes 0/0 to pooled
  figure; see \S4.4a \\[0.4em]
PLIF interaction recovery (pooled, rank1--10 oracle) --- \textbf{hybrid recomputation, valid} &
  \textbf{64/509 = 12.57\%} &
  same matched n=82 basis as above; per-complex mean 11.86\%, median 0\%; oracle headroom over the
  baseline row above is +3.34 points pooled (+3.05 points per-complex mean); see \S4.4a and
  \S4.5 \\[0.4em]
Confidence--PLIF recovery rank correlation --- original computation
  \textbf{[INVALIDATED, \S4.4a]} &
  \sout{Spearman rho = $-0.149$} (not significant) &
  n=78 (of the 80 PLIF-valid complexes in the original, now-superseded computation); computed
  against the wrong protein file --- see \S4.4a \\[0.4em]
Confidence--PLIF recovery rank correlation --- \textbf{hybrid recomputation, valid} &
  \textbf{Spearman rho = 0.046 (two-sided $p = 0.682$), not significant} &
  n=81 of the matched n=82 hybrid subset (one complex excluded, undefined recovery); independent
  average-rank Pearson calculation reproduced the same rho exactly; see \S4.4a \\
\bottomrule
\end{longtable}

\subsubsection*{4.4a Post-hoc invalidation of the original PLIF result, a valid hybrid
replacement, and the recomputed correlation (2026-07-21)}

The PLIF recovery rate and the confidence-vs-PLIF-recovery correlation originally reported for
this manuscript were computed by \path{eval_scripts/compute_plif.py} (this project's ``V1'' PLIF
pipeline) and were originally verified as \textit{reproducible} from \path{plif_recovery_results.csv}
and \path{master_eval_table_n122_with_plif.csv}. A subsequent code audit found that reproducibility
is not the same as methodological validity: V1 scores each predicted ligand pose's interaction
fingerprint against the \textbf{crystal protein file} (\path{protein_processed.pdb}) rather than
the \textbf{ESMFold-predicted protein file DiffDock-L actually used at inference time}
(\path{holo_aligned_predicted_protein.pdb}), confirmed by cross-checking V1's protein-file inputs
against \path{diffdock_dockgen_inputs.csv}. Every complex's PLIF fingerprint was therefore computed
in a structural context DiffDock-L never saw, which invalidates the resulting recovery numbers
regardless of how carefully the downstream pooling, denominator, or correlation arithmetic was
done.

\textbf{Invalidated.} Pooled PLIF interaction recovery 15.47\% (138/892, computed over 80 of the 91
RMSD-scored complexes; per-complex mean 17.3\%, median 0\%), and the confidence-vs-PLIF-recovery
Spearman rank correlation (rho = $-0.149$, n=78, $t(76) = -1.31$, not significant). Neither number
is cited as a finding anywhere in this draft; both are retained only as a transparent record of
what was measured and later retracted.

For denominator bookkeeping only (not a validity claim about the retracted V1 numbers): ProLIF
fingerprint computation itself failed outright on 11 of the 91 RMSD-scored complexes regardless of
protein file (7 from a numeric overflow in interaction-fingerprint indexing, \path{OverflowError:
Python integer ... out of bounds for uint32}; 3 from missing van der Waals radii for
metal-coordination atoms Fe/Mo; 1 from a residue-identity mismatch), leaving 80 as the denominator
for the retracted V1 pooled figure and 78 for the retracted V1 correlation (2 of the 80 lack a
\path{confidence_rank1} value, cause undiagnosed).

\textbf{Not affected.} Top-1 RMSD $<$ 2~\AA{} success rate (14/122 = 11.48\%) and ECE (0.203) are
computed by a separate pipeline (\path{build_master_table.py}, \path{compute_ece.py}) that does not
involve the ground-truth protein file used for PLIF scoring and is untouched by this bug.
Similarly, an exploratory (not significance-tested) rank correlation between confidence and binary
RMSD $<$ 2~\AA{} success (n=91 scored complexes) of approximately $+0.37$, in the expected
direction, does not depend on PLIF or the protein-file bug and is reported only for context, not
as a confirmed finding.

\textbf{Hybrid recomputation (complete).} A second, independently written pipeline (``V2'',
\path{dockgen_pilot_eval/}) uses the correct protein file but a broader ProLIF interaction-type
list (\path{interactions="all"}, 15 types, which double-counts $\pi$-stacking as both EdgeToFace
and FaceToFace) and reported pooled recovery 7.47\% (86/1151, matched n=82) --- never adopted as a
replacement for 15.47\%, since its interaction-list redundancy is itself a methodological problem.
The required combination --- V2's correct protein-file handling with V1's non-redundant 9-type
interaction list --- has now been run over the matched n=82 subset. Results:

\begin{itemize}
\item \textbf{Baseline (rank-1, confidence-selected pose):} pooled recovery 47/509 =
  \textbf{9.23\%}; per-complex mean 8.82\%, median 0\%. Source:
  \path{hybrid_plif_baseline_matched82.csv}, script \path{hybrid_plif_baseline.py}.
\item \textbf{Oracle (best of rank1--10 candidates per complex):} pooled recovery 64/509 =
  \textbf{12.57\%}; per-complex mean 11.86\%, median 0\%. Source:
  \path{hybrid_plif_oracle_matched82.csv}, script \path{hybrid_plif_oracle.py}.
\item \textbf{Headroom:} +3.34 percentage points pooled, +3.05 points per-complex mean (see \S4.5
  for the gate interpretation alongside the RMSD-side headroom).
\item \textbf{Denominator note:} one complex (\path{2wwc_1_CHT_2}) has zero true interactions
  under this 9-type list, making its per-complex recovery undefined (NaN); it contributes 0/0 to
  the pooled figures (no effect on the pooled percentages) and is excluded from the per-complex
  mean/median, which are therefore computed over n=81 non-null complexes.
\item Both files live under
  \path{~/docking_project/PoseBench/forks/DiffDock/results/dockgen_pilot_eval/} on the project's
  HPC environment.
\end{itemize}

These two numbers (9.23\% baseline, 12.57\% oracle) are this manuscript's reported PLIF recovery
figures, superseding the retracted 15.47\%.

\textbf{Confidence-vs-PLIF-recovery correlation (complete).} The correlation itself --- a distinct
computation from the recovery-rate recomputation above --- has also now been run against these
hybrid values. All 82 hybrid baseline rows matched a \path{confidence_rank1} value from
\path{master_eval_table_n122_with_plif.csv}; the one complex with zero true interactions and
therefore undefined recovery (\path{2wwc_1_CHT_2}) was excluded, leaving n=81. Spearman rho =
$0.046167845505$, two-sided $p = 0.682340974058$; an independent average-rank Pearson calculation
reproduced the same rho exactly. \textbf{The retracted rho = $-0.149$ (computed from the
invalidated V1 data) remains INVALIDATED and was not reused in any part of this computation.} As
with that earlier result, we read this only as a limited, non-causal statement: no statistically
significant monotonic relationship between DiffDock-L's confidence score and this corrected
PLIF-recovery rate was observed in this sample. Source: \path{eval_scripts/correlate_hybrid_plif.py},
\path{results/dockgen_pilot_eval/hybrid_plif_confidence_spearman.json}.

The RMSD-based success rate and ECE, the hybrid-recomputed PLIF recovery rate and its oracle
headroom, and the confidence-vs-PLIF-recovery correlation (rho = 0.046, $p = 0.682$, n=81, not
significant) are the measured Phase 1 numbers this draft can report as findings. Stratified
(pocket-similarity) reliability diagrams, the PDBBind in-distribution ECE anchor, AURC/risk--coverage,
and the base-DiffDock/FlowDock/classical-scoring comparators specified in \S3.1 have not
been computed. Consequently, this draft can report that DiffDock-L's confidence score is, in
absolute terms, poorly calibrated on this cross-docking benchmark, only weakly related to top-1
correctness, paired with a candidate pool whose chemical (PLIF) correctness is low and has only
modest headroom to improve by reranking, and shows no measured monotonic relationship to that
chemical correctness --- but it cannot yet report whether the observed miscalibration is
\textit{directional} (i.e., worse specifically under
distribution shift relative to an in-distribution anchor, as prior rank-correlation and
accuracy-degradation evidence would predict) or \textit{stratified} (concentrated in low
pocket-similarity complexes), because the comparison points needed to establish either claim have
not been measured in this project. We flag both as open, not as findings.

An earlier 65-complex pilot on an overlapping but not identical set reported success rate 18.5\%,
ECE 0.53, and PLIF recovery 6.0\%. Its computation scripts were not preserved, so this draft treats
those figures as unverifiable background only --- not as a comparison point --- and adopts the
n=122 result above as the project's sole reported baseline.

\subsection{Phase 2 status --- oracle-headroom gate and non-learned baselines (B1, B3) executed;
learned rescorer not yet trained}

The oracle-headroom gate (\S3.2, Step 0) --- the first required step before any rescoring model is
trained --- has now been run over per-rank RMSD and per-rank PLIF recovery for the retained
rank1--10 candidate pool, on the same 122-complex denominator as Phase 1.

\begin{table}[h]
\centering
\small
\begin{tabular}{@{}p{3.6cm}p{2.8cm}p{2.6cm}p{1.9cm}p{1.6cm}@{}}
\toprule
\textbf{Gate} & \textbf{Baseline (rank-1, confidence-selected)} & \textbf{Oracle (best of
rank1--10)} & \textbf{Headroom} & \textbf{Result vs. +5pp bar} \\
\midrule
RMSD $<$ 2~\AA{} & 14/122 = 11.48\% & 22/122 = 18.03\% & \textbf{+6.56 points} & \textbf{PASS} \\
PLIF recovery (pooled, hybrid-recomputed, \S4.4a) & 47/509 = 9.23\% & 64/509 = 12.57\% &
  \textbf{+3.34 points} & \textbf{FAIL} \\
\bottomrule
\end{tabular}
\end{table}

The two gates measure different things and disagree. RMSD headroom of +6.56 points means a
better-RMSD pose already exists, unchosen, in the confidence-ranked candidate pool for a
non-trivial number of complexes --- a \textit{selection} problem that a learned rescorer is
directly built to address. PLIF headroom of +3.34 points means that even an oracle that always
picked the single most chemically-correct candidate out of 10 barely improves over confidence's
own choice --- a \textit{sampling} limitation: the candidate pool itself rarely contains a much
better chemical match to select toward, so no reranking method, however good at selection, can
recover much there.

Applying the pre-registered decision rule per criterion, and treating either criterion clearing
the bar as sufficient to proceed (an effective logical OR across the two gates), \textbf{Phase 2
proceeds}, justified by the RMSD-side headroom alone. This draft does not extend that
justification to PLIF: no claim is made, and none should be inferred, that Interaction-Aware
Rescoring will materially improve PLIF-based chemical correctness --- the measured ceiling for
that improvement is +3.34 points regardless of rescorer quality. PLIF-derived features in the
proposed M1 model (\S3.2) are correspondingly framed as auxiliary signal for improving RMSD-based
selection, and any PLIF-recovery gain M1 shows in evaluation is to be reported as a secondary,
exploratory result rather than a headline contribution.

\textbf{Non-learned comparators B1 and B3 (executed, 2026-07-21).} Beyond the oracle gate (B2)
above, the two non-learned rerankers specified in \S3.2 have also been run on the same fixed
122-complex denominator, on the same retained rank1--10 candidate pool:

\begin{table}[h]
\centering
\small
\begin{tabular}{@{}p{3.4cm}p{2.4cm}p{2.6cm}p{6.0cm}@{}}
\toprule
\textbf{Comparator} & \textbf{Success rate} & \textbf{vs. B0 (14/122 = 11.48\%)} & \textbf{Basis} \\
\midrule
B0 --- raw confidence (rank-1) & 14/122 = 11.48\% & --- & existing baseline \\[0.3em]
B1 --- random-in-candidates (expected value) & 13.5667/122 = 11.12\% & $-0.36$ points & expected
  value from per-rank RMSD, 98/122 complexes evaluable; not an integer observation \\[0.3em]
B3 --- GNINA/Vina affinity rerank & 15/122 = 12.30\% & +0.82 points (+1 complex) & 111/122
  processed (SLURM array 12508307, 0 failures), 98/122 RMSD-evaluable; paired vs.\ B0: both-success
  11, B0-only 3, B3-only 4 \\
\bottomrule
\end{tabular}
\end{table}

B0, B1, and B3 all fall within an 11.12--12.30\% band on this benchmark. B3's one-complex gain over
B0 is not read as a clear improvement, given only 7 discordant pairs out of 122. This narrow spread
across a random baseline, the raw confidence baseline, and a classical affinity-based reranker
indicates that neither randomness nor classical scoring meaningfully exploits the RMSD-side
headroom identified by the oracle gate above. It motivates M1 as a learned alternative, but it
also sets a modest bar ($\sim$12\%) that M1 must clear by a defensible margin for Phase 2 to
represent more than a restatement of this band. M1 and M1-abl remain unexecuted; the single most
consequential open step, now that the gate and non-learned comparators are in hand, is the actual
training and evaluation of M1/M1-abl against B0/B1/B2/B3 as specified in \S3.2.

\subsection{Limitations}

(a) The Phase 1 protocol as designed includes five comparators and a stratified analysis; only one
comparator (DiffDock-L on DockGen-E, aggregate statistics) has been executed. Claims of
directional or stratified miscalibration are not yet supported by measurement. (b) 20 of 122
DockGen-E complexes (16\%) cannot be RMSD-scored under the current pipeline because of
non-standard ligand chemistry (peptide-like ligands, repeated glycine, metal ions) --- a
structural limitation of the DiffDock-L/PoseBench pipeline as deployed here, not resolved without
code-level changes, and a source of missing data rather than measured failure for those
complexes' correctness. (c) The per-complex PLIF-recovery distribution is heavily skewed (median
0\% in both the now-invalidated V1 computation and the valid hybrid recomputation, \S4.4a, for
baseline and oracle alike); Phase 2 will need to account for this in how it defines a PLIF-recovery
training target rather than treating it as a simple regression or balanced-classification signal.
(d) DockGen-E (n=122, 91 scored) gives low statistical power; any Phase 2 comparison will need the
PoseX-CD extension (n=1,312) before a significance claim is defensible, per the pre-registered
plan in \S3.2. (e) Comparator confidence scores (DiffDock-L's learned probability vs. Vina's
empirical scoring function vs. a trained rescorer's output) are not on a common semantic scale;
where a future draft reports ECE for multiple comparators, it must also report AURC/risk--coverage
as a semantics-free alternative, per the design in \S3.1. (f) The documentation discrepancy noted
in \S4.3 regarding the cause of the 11 runtime failures has been resolved in favor of the
execution-log breakdown (SVD ill-conditioning $\times$7, RDKit/exception-handling $\times$4); the
project's own public README still carries the outdated single-cause wording and needs a corrective
update outside this manuscript. (g) \textbf{(A historical denominator note about the now-retracted
V1 PLIF computation, \S4.4a; kept for the record, not for reuse.)} The retracted V1
pooled/mean/median recovery figures were computed over 80 of the 91 RMSD-scored complexes, not all
91, because ProLIF fingerprint computation itself failed outright on the remaining 11 (7
interaction-index overflow, 3 missing metal van der Waals radii, 1 residue-identity mismatch). The
valid hybrid recomputation (\S4.4a) uses a different denominator (matched n=82, one further
complex with an undefined/NaN recovery value) inherited from the V2 pipeline's matching procedure,
not this 80/91 breakdown; the two denominators should not be conflated. The same
fingerprint-computation failure modes remain a coverage risk for Phase 2 per-candidate feature
extraction (\S3.2). (h) \textbf{The confidence-vs-PLIF-recovery correlation has now been
recomputed against the valid hybrid PLIF values (\S4.4a).} The only earlier correlation computed
(n=78, rho = $-0.149$) used the now-retracted V1 PLIF pipeline (scored against the crystal, not
ESMFold-predicted, protein file) and remains uncited as a finding. The hybrid-value recomputation
(n=81, Spearman rho = $0.046167845505$, two-sided $p = 0.682340974058$, independently reproduced
via an average-rank Pearson calculation) finds no statistically significant monotonic
relationship --- a distinct analysis from the recovery-rate recomputation itself, and one that
should be read only as ``no clear rank relationship observed in this sample,'' not as evidence that
confidence and PLIF recovery are causally related or unrelated. (i) \textbf{The original PLIF
recovery measurement was invalidated by a ground-truth protein-file bug} (scored the predicted
ligand pose against the crystal protein file rather than the ESMFold-predicted protein file
DiffDock-L actually used); a hybrid recomputation against the correct protein file has since
produced valid replacement figures (9.23\% baseline, 12.57\% oracle, \S4.4a). The
confidence-vs-PLIF-recovery correlation, the one other PLIF-derived quantity from the original
computation, has also now been recomputed against the corrected data (rho = 0.046, n=81, not
significant; see (h) above). (j) \textbf{PDBBind training-subset (M1/M1-abl) exclusion bias.} Of
the 1,500-complex PDBBind subset drawn for M1/M1-abl training (\S3.2), after all retries and
code-level fixes 1,440/1,500 (96.0\%) complexes yielded valid outputs; the remaining 60 (4.0\%)
were excluded from the final training data: 37 SVD-ill-conditioned, 12 RDKit ligand-parsing
failures (overlapping a previously identified 236-complex list disproportionately containing
polysaccharide- and peptide-like ligands, consistent with PDBBind's own known $\sim$16\% RDKit
defect rate, arXiv:2411.01223), 6 TorchScript \path{radius.py} \texttt{RuntimeError} failures, 4
receptor-size-limit skips, and 1 CSV-field \texttt{AttributeError} (\S3.2). The direction of this
bias, not its magnitude, is the operative concern for the RDKit-related subset: M1 has
systematically fewer training examples on large, peptide-like ligands than a uniformly random 4\%
reduction would imply, and any accuracy comparison involving this chemotype should be read with
that caveat. \textbf{RMSD self-consistency for this final 1,440-complex sample has not yet been
recomputed} (\S3.2) --- this is a missing sanity check, not a null result, and should not be
inferred from the earlier, now-superseded 1,343-complex figure. (k) \textbf{PDBBind data-source substitution.}
M1/M1-abl's training structures come from a substitute Zenodo record (7014096) rather than the
source cited in DiffDock's own README (record 6408497, confirmed deleted; \S3.2), because the
original is no longer available. The substitute was prepared with a different
structure-preparation pipeline (Schr\"odinger Protein Prep Wizard vs. the original EquiBind
processing), so protonation states and coordinates may differ slightly from what DiffDock-L's
original authors used; this has not been empirically checked (e.g., via a backbone-RMSD spot check
between overlapping complex IDs) and is disclosed as an open reproducibility caveat, not a resolved
non-issue.

\section{Conclusion (interim)}

This draft reports a measured diagnosis, a measured oracle-headroom gate, and a specified but
not-yet-trained correction. On the DockGen-E cross-docking benchmark, DiffDock-L's top-1 pose is
correct (RMSD $<$ 2~\AA{}) for 14 of 122 complexes (11.48\%), and its confidence score has ECE
0.203 on the 91 complexes where correctness could be scored. An earlier reported PLIF recovery
figure (15.47\%, pooled) \textbf{is INVALIDATED (\S4.4a)}: it was computed against the wrong
(crystal, not ESMFold-predicted) protein file. A hybrid recomputation against the correct protein
file now gives a valid figure --- the top-1 pose recovers 9.23\% of reference protein--ligand
interactions pooled across the matched n=82 subset (47/509), against a rank1--10 oracle upper
bound of 12.57\% (64/509). A recomputed correlation between confidence and this corrected recovery
rate finds no statistically significant monotonic relationship (Spearman rho = 0.046, two-sided
$p = 0.682$, n=81) --- a limited, non-causal null result, not evidence that confidence and chemical
correctness are unrelated in any deeper sense.

The oracle-headroom gate (\S3.2 Step 0, \S4.5) that determines whether Phase 2 is worth building
has now been run, and it delivers two different verdicts depending on which correctness criterion
is used. On RMSD, headroom is substantial (14/122 to 22/122, +6.56 points): a better-RMSD pose
often already sits, unchosen, in DiffDock-L's own candidate pool, which is a selection failure a
learned rescorer can plausibly fix. On PLIF recovery, headroom is small (9.23\% to 12.57\%, +3.34
points): even a perfect oracle over the same 10 candidates recovers little more chemical
correctness than confidence already does, which is a sampling limitation no reranking method can
fix. We take the RMSD-side result alone as sufficient grounds to proceed to Phase 2, and we are
explicit that this draft does not claim, and will not claim without new evidence, that
Interaction-Aware Rescoring materially improves PLIF-based chemical correctness --- any such
improvement the trained model shows would be a secondary, exploratory finding, not the paper's
central contribution. Two non-learned rerankers run on the same candidate pool --- B1
(random-in-candidates, expected 11.12\%) and B3 (GNINA/Vina affinity reranking, 12.30\%) --- both
land within about a percentage point of confidence's own 11.48\% (\S3.2, \S4.5), indicating that
neither randomness nor classical scoring exploits the RMSD-side headroom on its own; this sets a
modest bar the proposed learned rescorer must clear. Whether the proposed rescoring model in fact
improves on the RMSD-based baseline by a margin beyond this band --- and whether any improvement
traces specifically to the PLIF-derived features rather than to cheaper signals such as ensemble
consistency --- is not yet known and remains the explicit subject of the pre-registered,
not-yet-executed training and evaluation protocol in \S3.2.

\section*{Code and Data Availability}

Code, configuration, and result tables for the measured Phase 1 results are available at
\path{diffdock-interaction-rescoring}. This draft's \S4.3--4.4 numbers match the CSVs and
execution log in the repository, but not the repository README's failure-cause narrative: as of
this draft, the README still attributes all 11 runtime failures to a single group-aggregation
bug, which \S4.3 and \S4.6(f) above identify as incorrect (the primary execution log shows 7 SVD
ill-conditioning failures and 4 RDKit/exception-handling failures instead). This is an
outdated-documentation issue in the public repository, not a discrepancy in the numbers reported
here, and has not yet been corrected outside this manuscript. Result files, including the main
evaluation table (\path{master_eval_table_n122_with_plif.csv}), ECE/reliability-diagram data
(\path{ece_reliability_bins_n91.csv}, \path{reliability_diagram_n91.png}), and PLIF recovery
results (\path{plif_recovery_results.csv}), are under \path{results/dockgen_pilot/}, with
reproduction scripts under \path{results/dockgen_pilot/eval_scripts/}. \textbf{\path{plif_recovery_results.csv}
and the PLIF columns of \path{master_eval_table_n122_with_plif.csv} are invalidated (\S4.4a):}
\path{compute_plif.py} scores the predicted ligand pose against the crystal protein file rather
than the ESMFold-predicted protein file DiffDock-L actually used at inference time. These files
are retained in the repository as a record, not as a source of valid PLIF numbers. \textbf{The
valid replacement PLIF figures reported in \S4.4/\S4.4a (9.23\% baseline, 12.57\% oracle) come
from a separate hybrid pipeline}: \path{hybrid_plif_baseline_matched82.csv} and
\path{hybrid_plif_oracle_matched82.csv} (scripts \path{hybrid_plif_baseline.py},
\path{hybrid_plif_oracle.py}), under \path{results/dockgen_pilot_eval/} --- not yet
cross-referenced against the public repository's directory layout at the time of this draft;
confirm the path before final submission. \textbf{The confidence-vs-PLIF-recovery correlation
against these hybrid values has since been recomputed (\S4.4a, \S4.6(h)):}
\path{results/dockgen_pilot_eval/hybrid_plif_confidence_spearman.json}, script
\path{eval_scripts/correlate_hybrid_plif.py}. Reproducing these numbers from scratch currently
requires manual work beyond cloning the repository: \path{eval_scripts/*.py} hardcode absolute
paths from the project's HPC compute environment and will not run unmodified elsewhere, and two of
the three scripts require input files not yet tracked in the repository ---
\path{build_master_table.py} needs \path{diffdock_dockgen_inputs.csv}, and \path{compute_plif.py}
needs \path{data/dockgen_set/}. The tracked CSV outputs in \path{results/dockgen_pilot/} can be
used directly without rerunning these scripts. DockGen-E is redistributed by PoseBench
(BioinfoMachineLearning/PoseBench) and originates from Corso et al. (arXiv:2402.18396). The
oracle-headroom gate and the two non-learned comparators B1 and B3 (\S4.5) have been executed ---
result files \path{random_baseline_rmsd.json} and \path{gnina_b3_selected.csv}/
\path{gnina_b3_summary.json} under \path{results/dockgen_pilot_eval/} --- but no Phase 2 learned
rescoring model (M1/M1-abl) code or trained artifact exists yet; this section will be updated once
that model is trained and evaluated. M1's training data (a 1,500-complex PDBBind subset) is being
provisioned from a substitute Zenodo source (record 7014096) because the source cited in
DiffDock's own README (record 6408497) was found deleted; this substitution, and its unverified
structure-preparation difference, are disclosed in \S3.2 and must be carried into any future
release of the M1 training data itself.

% Bibliography: see refs.bib (split out from this file's former inline
% thebibliography block on 2026-07-21). Same citation keys as before.
% Inline citations in the running text above are still plain
% "(Author et al., Year)" text, not \cite{}; the citation-key pass into
% the body has not been done (see wiki/meta/manuscript-status.md, open
% issue 4).
\bibliographystyle{plain}
\bibliography{refs}

\end{document}
